\documentclass[11pt,a4paper]{article}

% ─── Encoding & language ───────────────────────────────────────────────────
\usepackage[utf8]{inputenc}
\usepackage[T1]{fontenc}
\usepackage[spanish,es-tabla]{babel}
\usepackage{lmodern}

% ─── Page geometry ─────────────────────────────────────────────────────────
\usepackage[a4paper,margin=2.3cm,headheight=15pt]{geometry}

% ─── Colors ────────────────────────────────────────────────────────────────
\usepackage{xcolor}
\definecolor{apexorange}{HTML}{F97316}
\definecolor{apexsteel}{HTML}{1F2937}
\definecolor{apexgray}{HTML}{6B7280}
\definecolor{apexlight}{HTML}{F3F4F6}
\definecolor{apexgreen}{HTML}{10B981}
\definecolor{apexred}{HTML}{EF4444}
\definecolor{apexblue}{HTML}{3B82F6}
\definecolor{apexamber}{HTML}{F59E0B}

% ─── Typography ────────────────────────────────────────────────────────────
\usepackage{titlesec}
\titleformat{\section}
  {\Large\bfseries\color{apexorange}}
  {\thesection.}{0.6em}{}
  [\vspace{-0.4em}\color{apexorange}\rule{\linewidth}{0.6pt}]
\titleformat{\subsection}
  {\large\bfseries\color{apexsteel}}
  {\thesubsection}{0.5em}{}
\titleformat{\subsubsection}
  {\normalsize\bfseries\color{apexgray}}
  {\thesubsubsection}{0.4em}{}

\titlespacing*{\section}{0pt}{1.4em}{0.6em}
\titlespacing*{\subsection}{0pt}{1em}{0.4em}

% ─── Header / footer ───────────────────────────────────────────────────────
\usepackage{fancyhdr}
\pagestyle{fancy}
\fancyhf{}
\renewcommand{\headrulewidth}{0.4pt}
\renewcommand{\headrule}{\hbox to\headwidth{\color{apexorange}\leaders\hrule height \headrulewidth\hfill}}
\fancyhead[L]{\textcolor{apexorange}{\textbf{APEX}} \textcolor{apexgray}{\small Manual de Uso}}
\fancyhead[R]{\textcolor{apexgray}{\small\thepage}}

% ─── Lists ─────────────────────────────────────────────────────────────────
\usepackage{enumitem}
\setlist[itemize]{leftmargin=1.3em,itemsep=0.2em,topsep=0.3em}
\setlist[enumerate]{leftmargin=1.5em,itemsep=0.3em,topsep=0.4em}

% ─── Tables ────────────────────────────────────────────────────────────────
\usepackage{booktabs}
\usepackage{array}
\usepackage{tabularx}
\renewcommand{\arraystretch}{1.25}
\newcolumntype{L}[1]{>{\raggedright\arraybackslash}p{#1}}

% ─── Boxes (callouts) ──────────────────────────────────────────────────────
\usepackage[most]{tcolorbox}

\newtcolorbox{infobox}[1][]{
  colback=apexblue!7,colframe=apexblue,
  fonttitle=\bfseries,coltitle=white,
  title={\textcolor{white}{Nota}},
  arc=2pt,boxrule=0.6pt,left=8pt,right=8pt,top=6pt,bottom=6pt,
  #1
}
\newtcolorbox{warnbox}[1][]{
  colback=apexamber!10,colframe=apexamber,
  fonttitle=\bfseries,coltitle=white,
  title={\textcolor{white}{Importante}},
  arc=2pt,boxrule=0.6pt,left=8pt,right=8pt,top=6pt,bottom=6pt,
  #1
}
\newtcolorbox{dangerbox}[1][]{
  colback=apexred!8,colframe=apexred,
  fonttitle=\bfseries,coltitle=white,
  title={\textcolor{white}{Cuidado}},
  arc=2pt,boxrule=0.6pt,left=8pt,right=8pt,top=6pt,bottom=6pt,
  #1
}
\newtcolorbox{tipbox}[1][]{
  colback=apexgreen!8,colframe=apexgreen,
  fonttitle=\bfseries,coltitle=white,
  title={\textcolor{white}{Tip}},
  arc=2pt,boxrule=0.6pt,left=8pt,right=8pt,top=6pt,bottom=6pt,
  #1
}

% ─── Inline code ───────────────────────────────────────────────────────────
\usepackage{xparse}
\NewDocumentCommand{\code}{m}{\colorbox{apexlight}{\texttt{\small #1}}}
\NewDocumentCommand{\btn}{m}{\textbf{\textcolor{apexorange}{[#1]}}}
\NewDocumentCommand{\menu}{m}{\textbf{\textcolor{apexsteel}{#1}}}

% ─── Hyperlinks ────────────────────────────────────────────────────────────
\usepackage[colorlinks=true,linkcolor=apexorange,urlcolor=apexorange,citecolor=apexorange]{hyperref}

% ─── Title metadata ────────────────────────────────────────────────────────
\title{}
\author{}
\date{}

% ═══════════════════════════════════════════════════════════════════════════
\begin{document}

% ─── Cover page ────────────────────────────────────────────────────────────
\begin{titlepage}
  \begin{center}
    \vspace*{2cm}
    {\Huge\bfseries\textcolor{apexorange}{APEX}}\\[0.3cm]
    {\Large\bfseries\textcolor{apexsteel}{CONCRETE}}\\[2cm]

    \rule{0.6\linewidth}{1pt}\\[0.5cm]
    {\LARGE\bfseries Manual de Uso}\\[0.3cm]
    {\large\textcolor{apexgray}{Sistema interno de facturas y planilla}}\\[0.3cm]
    \rule{0.6\linewidth}{1pt}\\[3cm]

    \begin{tabular}{rl}
      \textcolor{apexgray}{\textbf{Versión:}}    & 1.0 \\
      \textcolor{apexgray}{\textbf{Frontend:}}   & \url{https://apex-tau-ebon.vercel.app} \\
      \textcolor{apexgray}{\textbf{Backend:}}    & \url{https://apex-w9h5.onrender.com} \\
      \textcolor{apexgray}{\textbf{Actualizado:}}& Abril 2026 \\
    \end{tabular}

    \vfill
    \textcolor{apexgray}{\small Documento interno --- Apex Concrete Inc.}
  \end{center}
\end{titlepage}

% ─── Table of contents ─────────────────────────────────────────────────────
\renewcommand*\contentsname{\textcolor{apexorange}{Contenido}}
\tableofcontents
\thispagestyle{fancy}
\newpage

% ═══════════════════════════════════════════════════════════════════════════
\section{Introducción}

Apex es la aplicación interna que reemplaza el flujo manual de QuickBooks Online (QBO) y el cálculo de planilla en hojas de cálculo. Sincroniza facturas desde QBO, calcula automáticamente el pago al colaborador según el tipo de trabajo, y genera reportes exportables a Excel y PDF.

\subsection{Componentes principales}

\begin{itemize}
  \item \menu{Facturas}: lista y detalle de cada factura sincronizada de QBO.
  \item \menu{Reportes $\rightarrow$ Salarios}: planilla por colaborador y rango de fechas.
  \item \menu{Reportes $\rightarrow$ EPOs}: facturas de tipo EPO que requieren ajuste manual.
  \item \menu{Configuración $\rightarrow$ Colaboradores}: registro maestro de quien recibe pagos.
  \item \menu{Configuración $\rightarrow$ Conexiones}: vinculación con QBO y sincronización.
\end{itemize}

\begin{infobox}
La aplicación corre en infraestructura gratuita (Render). El backend puede tardar hasta 15 segundos en responder si llevaba inactivo varias horas. Es normal --- el primer request lo despierta.
\end{infobox}

% ═══════════════════════════════════════════════════════════════════════════
\section{Setup inicial (una sola vez)}

\subsection{Registrar colaboradores}

Antes de cualquier sincronización, hay que registrar a cada colaborador en la app. Los nombres aquí son la única referencia que tiene Apex para asignar pagos.

\begin{enumerate}
  \item Ir a \menu{Configuración $\rightarrow$ Colaboradores}.
  \item Click en \btn{+ Agregar}.
  \item Llenar:
    \begin{itemize}
      \item \textbf{Nombre:} debe coincidir EXACTAMENTE con lo que Emily escribirá en QBO.
      \item \textbf{Email:} opcional, usado para envío de recibos.
      \item \textbf{Color:} para identificación visual.
    \end{itemize}
  \item Guardar.
\end{enumerate}

\begin{tipbox}
Usar nombres únicos y sin ambigüedad. Por ejemplo, en lugar de \code{ED} (que podría confundirse con \textit{REPLACED}), usar \code{Eduardo}. Nombres de dos palabras (\code{Juan Perez}) son más seguros que solo \code{Juan}.
\end{tipbox}

\subsection{Conectar QuickBooks}

\begin{enumerate}
  \item \menu{Configuración $\rightarrow$ Conexiones}.
  \item Click en \btn{Conectar con QuickBooks}.
  \item Iniciar sesión con la cuenta de QBO de Apex.
  \item Autorizar acceso.
\end{enumerate}

La conexión se mantiene activa. Solo hay que reconectar si QBO la revoca o si vence el token (raro).

% ═══════════════════════════════════════════════════════════════════════════
\section{Cómo se calculan los pagos}

\subsection{Tabla de tarifas automáticas}

Apex lee los \textit{line items} de cada factura. El producto/servicio determina si se paga y cuánto:

\begin{center}
\begin{tabularx}{\linewidth}{L{4.5cm} X r}
\toprule
\textbf{Tipo de trabajo (en QBO)} & \textbf{Cómo se paga} & \textbf{Tarifa} \\
\midrule
Pour Mono Slab / Mono Slab & Automático & \$1.00 / SF \\
Sidewalk                   & Automático & \$0.75 / SF \\
City Sidewalk              & Automático & \$0.75 / SF \\
Extra Sidewalk             & Automático & \$0.75 / SF \\
Drives \& Walks            & Automático & \$0.75 / SF \\
Driveways                  & Automático & \$0.75 / SF \\
\midrule
MATERIAL                   & Nunca se paga (solo se muestra) & \$0 \\
EPO: \textit{<trabajo>}    & Manual, vía pestaña EPOs & \$0 auto \\
Turn Down, Undercut, Muck Out, Replacement, Patch, otros & Manual, en modal de factura & \$0 auto \\
\bottomrule
\end{tabularx}
\end{center}

\subsection{Asignación del colaborador}

Apex lee el campo \textbf{``Memo on statement (hidden)''} de cada factura en QBO. Busca dentro de ese texto el nombre de un colaborador registrado.

\begin{center}
\begin{tabularx}{\linewidth}{L{4cm} L{4cm} X}
\toprule
\textbf{Memo en QBO} & \textbf{Colaborador en app} & \textbf{Resultado} \\
\midrule
\code{Juan Perez}    & \code{Juan}          & Asigna a Juan \\
\code{JUAN, MATERIAL}& \code{Juan}          & Asigna a Juan \\
\code{REPLACED}      & \code{Ed}            & \textbf{NO asigna} (palabra parcial) \\
\code{JUANITO}       & \code{Juan}          & \textbf{NO asigna} (no es palabra completa) \\
\code{CITYWALK}      & cualquiera           & \textbf{NO asigna} (sin coincidencia) \\
\code{Juan Perez}    & \code{Juan} y \code{Juan Perez} ambos existen & Asigna al más específico (\code{Juan Perez}) \\
\bottomrule
\end{tabularx}
\end{center}

\begin{warnbox}
Si el memo no contiene el nombre de un colaborador registrado, la factura queda en estado \textbf{``Sin asignar''}. Esto es comportamiento intencional: prefiere quedar sin asignar que asignar al colaborador equivocado.
\end{warnbox}

% ═══════════════════════════════════════════════════════════════════════════
\section{Flujo de Emily en QBO}

Para que la app funcione bien, Emily debe seguir estas reglas al crear o editar facturas en QuickBooks:

\subsection{Memo on statement (hidden)}

Escribir el nombre del colaborador exactamente como está registrado en Apex. Puede ir solo o acompañado de otra info, separado por espacio o coma.

\begin{itemize}
  \item Correcto: \code{Juan Perez}
  \item Correcto: \code{Juan Perez, EXTRA}
  \item Incorrecto: \code{Juanito} (no matchea con \code{Juan})
  \item Incorrecto: dejar vacío --- la factura quedará sin asignar
\end{itemize}

\subsection{Line items (Product/Service)}

Cada línea de la factura debe usar uno de los \textit{products} estándar listados en la tabla de tarifas. La cantidad (Qty) en pies cuadrados.

\begin{itemize}
  \item Si es un \textbf{EPO}, el formato debe ser \code{EPO: <nombre del trabajo>}.
  \item Si es un trabajo no estándar (Replacement, Patch, etc.), Apex lo dejará en \$0 y Emily debe ingresar el pago manualmente desde el modal de la factura en la app.
\end{itemize}

% ═══════════════════════════════════════════════════════════════════════════
\section{Sincronización}

Apex no se conecta a QBO en tiempo real. Hay que sincronizar manualmente o automáticamente.

\subsection{Tipos de sync}

\begin{center}
\begin{tabularx}{\linewidth}{L{3.5cm} X L{2.5cm}}
\toprule
\textbf{Tipo} & \textbf{Para qué sirve} & \textbf{Velocidad} \\
\midrule
\btn{Sync rápido (1 sem)} & Día a día. Trae solo los últimos 7 días. & $\sim$10 segundos \\
\btn{Sync completo}       & Cuando se agrega un colaborador nuevo, se cambia una regla, o se sospecha de datos viejos incorrectos. Trae todas las facturas. & 1--3 minutos \\
\bottomrule
\end{tabularx}
\end{center}

\subsection{Cuándo correr cada uno}

\begin{itemize}
  \item \textbf{Diariamente o antes de revisar:} Sync rápido.
  \item \textbf{Después de crear un colaborador nuevo:} Sync completo (para que reasigne facturas viejas que ahora matchean).
  \item \textbf{Después de cambiar tarifas:} Sync completo (recalcula pagos).
  \item \textbf{Antes de exportar planilla:} Sync rápido para asegurar datos al día.
\end{itemize}

\begin{infobox}
La sincronización \textbf{nunca} pisa los pagos manuales. Si Emily editó el pago de una factura desde la app, ese valor sobrevive cualquier sync futuro.
\end{infobox}

% ═══════════════════════════════════════════════════════════════════════════
\section{Revisión semanal recomendada}

Cada lunes (o antes de cerrar la semana de pago):

\subsection{Paso 1: Sincronizar}
Click en \btn{Sync rápido (1 sem)} desde \menu{Configuración}.

\subsection{Paso 2: Facturas sin asignar}
\begin{enumerate}
  \item Ir a \menu{Facturas}.
  \item Filtrar por colaborador = \textit{Sin asignar}.
  \item Para cada una, abrir el modal y asignar manualmente al colaborador correcto.
\end{enumerate}

\subsection{Paso 3: Revisar pagos automáticos}
\begin{enumerate}
  \item En \menu{Facturas}, revisar facturas recientes.
  \item Si el pago calculado no coincide con lo acordado (ej. Emily pactó tarifa especial), abrir modal y editar el campo \textbf{``Pago colaborador''}.
  \item Aparecerá badge \btn{MANUAL} indicando que es override y no se recalculará.
\end{enumerate}

\subsection{Paso 4: EPOs}
\begin{enumerate}
  \item \menu{Reportes $\rightarrow$ EPOs}.
  \item Cada EPO requiere SF (cantidad) o pago manual ingresado por Emily.
  \item Una vez ingresado, se cuenta automáticamente en la planilla del colaborador.
\end{enumerate}

\subsection{Paso 5: Pagos off-invoice}
\begin{enumerate}
  \item \menu{Reportes $\rightarrow$ Salarios $\rightarrow$ Agregar pago manual}.
  \item Usar para bonus, adelantos, ajustes que no están ligados a una factura.
\end{enumerate}

% ═══════════════════════════════════════════════════════════════════════════
\section{Pago de planilla}

\subsection{Generar reporte}

\begin{enumerate}
  \item \menu{Reportes $\rightarrow$ Salarios}.
  \item Seleccionar rango de fechas (típicamente lunes a domingo de la semana cerrada).
  \item Opcional: filtrar por colaborador específico.
  \item Revisar tabla en pantalla --- muestra cada factura, EPO modificado, y pago manual contribuyendo al total.
\end{enumerate}

\subsection{Exportar}

\begin{itemize}
  \item \btn{Exportar Excel}: detalle por colaborador, una hoja por persona.
  \item \btn{Exportar PDF}: resumen visual para enviar por email o imprimir.
  \item \btn{Enviar email}: si el colaborador tiene email registrado, manda PDF directo.
\end{itemize}

% ═══════════════════════════════════════════════════════════════════════════
\section{Mantenimiento de colaboradores}

\subsection{Limpiar basura}

Si por algún motivo se acumularon registros de colaboradores que no deberían existir:

\begin{enumerate}
  \item \menu{Configuración $\rightarrow$ Colaboradores}.
  \item Click en \btn{Limpiar sin facturas}: borra automáticamente todos los que no tienen ninguna factura ni pago manual asignado.
\end{enumerate}

\subsection{Eliminar colaboradores con facturas}

Si un colaborador tiene facturas pero igual hay que removerlo (ej. duplicado, error de nombre):

\begin{enumerate}
  \item Marcar el checkbox del colaborador.
  \item Click en \btn{Eliminar permanentemente}.
  \item Confirmar el aviso. Si tiene facturas, confirmar el segundo aviso de \textit{forzar}.
  \item Las facturas quedan en estado ``Sin asignar''. Reasignar manualmente desde el modal de cada factura.
\end{enumerate}

\begin{dangerbox}
La eliminación permanente NO se puede deshacer. Si solo querés desactivarlo (que no aparezca en listas pero conserve sus pagos históricos), usar \btn{Desactivar} en su lugar.
\end{dangerbox}

% ═══════════════════════════════════════════════════════════════════════════
\section{Reglas de oro}

\begin{itemize}
  \item El \textbf{pago manual siempre gana} sobre el cálculo automático. Sync nunca lo pisa.
  \item Si el memo en QBO no matchea ningún colaborador, la factura queda \textbf{Sin asignar} (no se asigna mal por error).
  \item EPOs y trabajos no estándar (Replacement, Patch, Turn Down, etc.) \textbf{nunca pagan automáticamente}; siempre requieren override manual.
  \item Cambios en tarifas o tipos de trabajo solo aplican después de un sync completo.
  \item \textbf{No} eliminar colaboradores recién creados sin antes correr sync completo: la asignación de facturas viejas necesita el sync para detectarlos.
\end{itemize}

% ═══════════════════════════════════════════════════════════════════════════
\section{Solución de problemas}

\subsection{Una factura quedó ``Sin asignar''}

\textbf{Causa:} El memo en QBO no contiene el nombre del colaborador como palabra completa, o el colaborador no está registrado en la app.

\textbf{Solución:}
\begin{enumerate}
  \item Verificar que el colaborador esté en \menu{Configuración $\rightarrow$ Colaboradores}.
  \item Abrir la factura en \menu{Facturas} y asignar manualmente.
  \item Para evitarlo en futuras facturas, corregir el memo en QBO.
\end{enumerate}

\subsection{El pago calculado no es correcto}

\textbf{Causa:} El \textit{product/service} en QBO no matchea ninguno de los tipos auto-pagables, o la cantidad (Qty) está mal.

\textbf{Solución:}
\begin{enumerate}
  \item Verificar el line item en QBO.
  \item Si el tipo no es auto-pagable (ej. \textit{Replacement}), abrir modal de la factura en la app y editar el pago manualmente.
\end{enumerate}

\subsection{No aparece una factura recién creada en QBO}

\textbf{Causa:} Aún no se sincronizó.

\textbf{Solución:} Correr \btn{Sync rápido (1 sem)} desde \menu{Configuración}.

\subsection{El backend tarda mucho en responder}

\textbf{Causa:} Render free tier suspende el servicio tras 15 minutos sin uso. El primer request lo despierta.

\textbf{Solución:} Esperar hasta 15 segundos en el primer click. Los siguientes serán inmediatos.

\subsection{EPOs no aparecen en la pestaña EPOs}

\textbf{Causa:} El \textit{product/service} en QBO no empieza con \code{EPO:}.

\textbf{Solución:} Verificar formato en QBO. Debe ser exactamente \code{EPO: <nombre del trabajo>}.

% ═══════════════════════════════════════════════════════════════════════════
\section{Resumen visual del flujo}

\begin{center}
\begin{tcolorbox}[colback=apexlight,colframe=apexorange,boxrule=1pt,arc=4pt]
\textbf{1.} Emily crea factura en QBO con memo = nombre colab + line items correctos. \\[0.4em]
\textbf{2.} Apex sincroniza (rápido o completo). \\[0.4em]
\textbf{3.} App calcula pago según tipo de trabajo y asigna colaborador. \\[0.4em]
\textbf{4.} Emily revisa en app: ajusta sin asignar, EPOs, pagos manuales. \\[0.4em]
\textbf{5.} Generar reporte de planilla por rango de fechas. \\[0.4em]
\textbf{6.} Exportar a Excel o PDF y enviar al colaborador.
\end{tcolorbox}
\end{center}

\vfill

\begin{center}
\textcolor{apexgray}{\small \textit{Documento interno de Apex Concrete --- Para soporte técnico contactar al desarrollador.}}
\end{center}

\end{document}
